\documentclass[11pt,a4paper]{article}
% Shared preamble for the three companion documents.  Deliberately minimal:
% only packages a bare TeX Live carries, since every figure is built from
% LaTeX `picture` primitives rather than from TikZ or an external plotter.
\usepackage[utf8]{inputenc}
\usepackage[T1]{fontenc}
\usepackage{amsmath,amssymb,amsthm}
\usepackage{booktabs}
\usepackage{algorithm}
\usepackage{algpseudocode}
\usepackage{xcolor}
\usepackage[margin=2.4cm]{geometry}
\usepackage[hidelinks]{hyperref}

\newcommand{\Z}{Z}
\newcommand{\Dset}{D}
\newcommand{\E}{E(P_{\Dset})}
\newcommand{\R}{\mathbb{R}}
\newcommand{\Zint}{\mathbb{Z}}
\newcommand{\fig}[1]{\input{figs/#1}}

% the legend shared by every criterion-space figure
\newcommand{\figlegend}{%
  \par\smallskip\noindent{\footnotesize
  \textcolor{green!45!black}{$\blacksquare$}~efficient point\quad
  \textcolor{black!45}{$\square$}~dominated point\quad
  \textcolor{blue!60!black}{$\square$}~the box being settled\quad
  \textcolor{red!35}{$\blacksquare$}~what the cut removes\quad
  \textcolor{red!65!black}{$\square$}~the children left\quad
  \textcolor{red!65!black}{$\circ$}~$x_B$\quad
  \textcolor{orange!85!black}{$\circ$}~the centre cut on\par}}

\newtheorem{proposition}{Proposition}
\newtheorem{remark}{Remark}


\title{Results, with the charts behind them\\
\large what was measured, in which unit, and what the unit hides}
\author{}
\date{}

\begin{document}
\maketitle

\begin{abstract}
Every number here comes from a run of the shipped code in exact rational
arithmetic, cross-checked against an independent reference, and every chart is
generated from those numbers by \texttt{docs/figs/regenerate.py}.  The losses
are reported at the same length as the wins, and so is the one finding that
outlasts all of them: on this machine wall-clock time cannot resolve the
effects being measured, so most of what follows is counted in \emph{work} ---
sub-problems, branch-and-bound nodes, simplex pivots --- and not in seconds.
\end{abstract}

\tableofcontents

\section{What is being compared}

Two exact methods answer $(P_E)$ and agree on every instance.  They differ in
where they keep the unexplored part of the problem, and that decides the cost:

\begin{description}
\item[decision space] truncates $\Dset$ with Sylva--Crema cuts.  Each cut is a
  disjunction, so it costs $p$ binaries and $p+1$ big-$M$ rows, and every
  sub-problem is strictly larger than the last.
\item[criterion space] keeps a list of boxes in $\R^p$.  The same truncation is
  a box subtraction; every sub-problem is $\Dset$ plus a few ordinary rows, and
  \textbf{the model never grows}.
\end{description}

\section{The margin, and how it moves with $p$}

Thirty instances per row, $n = 6$.  The statistic is the \textbf{median of the
per-instance ratios}: a ratio of totals would be decided by the single heaviest
instance, which is the one thing a small sample gets most wrong.

\begin{table}[h]
\centering
\caption{Decision space against criterion space, thirty instances per row.}
\label{tab:cmargin}
\begin{tabular}{crrrrr}
\toprule
$p$ & decision & criterion & ratio (median) & IQR & max \\
\midrule
2 & $0.07$ s & $0.04$ s & $1.52\times$  & $[1.25,\ 1.91]$   & $6.0\times$ \\
3 & $0.34$ s & $0.12$ s & $2.69\times$  & $[1.87,\ 3.72]$   & $10.1\times$ \\
4 & $0.50$ s & $0.13$ s & $4.99\times$  & $[2.39,\ 11.62]$  & $76.3\times$ \\
5 & $2.19$ s & $0.25$ s & $7.93\times$  & $[3.61,\ 19.98]$  & $33.8\times$ \\
6 & $3.67$ s & $0.31$ s & $14.57\times$ & $[7.90,\ 25.17]$  & $65.7\times$ \\
\bottomrule
\end{tabular}
\end{table}

\begin{center}
\fig{r-campaign-margin}
\end{center}

\noindent
Three of the 150 runs were stopped by a $45$ s budget --- one at $p=5$, two at
$p=6$ --- and are excluded rather than counted as their budget.  That censoring
is \emph{one-sided}: it can only make the decision-space method look better, so
every margin here is conservative.  Every instance that finished gave the same
optimum on both sides.

\paragraph{An earlier table said the curve turned over.  It does not.}  Four
instances per $p$ gave $1.74$, $2.99$, $6.92$, $103$ and $24.3$, and the text
went on to explain the turn between $p=5$ and $p=6$.  There was nothing to
explain: that $103\times$ was one instance supplying $110.95$ of $112.77$
seconds.  With thirty instances the progression is clean and monotone.  The old
figures were not wrong as measurements --- they are what those four instances
did --- but they were quoted as a trend, and four runs cannot carry one.

\section{Where the margin comes from, counted}

Wall-clock cannot resolve a small effect on this machine, so the campaign
counts the deterministic work as well: sub-programs solved and branch \& bound
nodes explored, identical on every run and every machine.

\begin{table}[h]
\centering
\caption{Ratios of decision space over criterion space, same thirty instances
per row.}
\begin{tabular}{crrr}
\toprule
$p$ & sub-programs & b\&b nodes & time \\
\midrule
2 & $1.72\times$ & $1.28\times$ & $1.52\times$ \\
3 & $1.40\times$ & $1.79\times$ & $2.69\times$ \\
4 & $1.19\times$ & $2.68\times$ & $4.99\times$ \\
5 & $1.17\times$ & $3.16\times$ & $7.93\times$ \\
6 & $\mathbf{1.02\times}$ & $\mathbf{5.37\times}$ & $\mathbf{14.57\times}$ \\
\bottomrule
\end{tabular}
\end{table}

\begin{center}
\fig{r-campaign-work}\\[2pt]
\fig{r-campaign-work-legend}
\end{center}

\noindent
\textbf{This is the whole cost model, read off the measurement.}  At $p=6$ the
two methods solve \emph{the same number of sub-programs} --- $1.02\times$, down
steadily from $1.72\times$ at $p=2$ --- and one still takes $14.57\times$ as
long.  The entire margin has moved into the cost of \emph{one} sub-problem:
the branch \& bound inside it explores $5.37\times$ the nodes, because by then
it carries $p$ binaries and $p+1$ big-$M$ rows per cut while the box search
carries none.  The branching that was supposed to work against the box list
wins the sub-program count at small $p$ and gives it back as $p$ grows; the
other term compounds.

\section{Is the margin real, or is the baseline handicapped?}

A $103\times$ margin --- and a $204\times$ one on front enumeration --- invites
the suspicion that the loser was crippled rather than beaten, so the losing
method was profiled instead of merely quoted.

\paragraph{Where the time goes.}  On the instance that takes $137$ s to produce
a 30-vector front in decision space:

\begin{table}[h]
\centering
\caption{Profile of the decision-space enumeration, $137.38$ s in total.}
\begin{tabular}{lrrr}
\toprule
 & time & share & calls \\
\midrule
probe (the integer program) & $137.21$\,s & $\mathbf{100\%}$ & $31$ \\
repair to an efficient point & $0.09$\,s & $0\%$ & $30$ \\
building and adding the cut & $0.02$\,s & $0\%$ & $30$ \\
\bottomrule
\end{tabular}
\end{table}

\noindent
Over those 31 calls, branch \& bound explores $14983$ nodes, and the model
grows from $5$ columns and $7$ rows to $95$ and $127$ --- exactly $p$ binaries
and $p+1$ rows per cut.  Nothing outside the sub-problems is worth optimising,
and the growth of the sub-problems is precisely what the criterion-space method
does not have.  \textbf{The cost is structural.}

\paragraph{The one candidate handicap, tested.}  The probe only needs
\emph{some} feasible point of the truncated region, yet it maximises a linear
direction.  Dropping the objective looks strictly cheaper --- stop at the first
integer point found.  Measured, it is four to thirteen times \emph{worse}:

\begin{table}[h]
\centering
\caption{Probing by maximising a direction against probing with a zero
objective.  Same instances, same cuts.}
\begin{tabular}{lrrr}
\toprule
 & cuts & b\&b nodes & time \\
\midrule
$n=4$, maximise a direction & $9$ & $658$  & $\mathbf{1.35}$\,s \\
$n=4$, zero objective       & $9$ & $1397$ & $5.33$\,s \\
$n=5$, maximise a direction & $9$ & $372$  & $\mathbf{0.79}$\,s \\
$n=5$, zero objective       & $9$ & $2105$ & $10.15$\,s \\
\bottomrule
\end{tabular}
\end{table}

\begin{center}
\fig{r-probe-cuts}\\[6pt]
\fig{r-probe-nodes}\\[2pt]
\fig{r-probe-legend}
\end{center}

\noindent
The two charts are the whole finding: the top one is flat and the bottom one is
not.  \textbf{The cut counts are identical} --- the direction changes neither
how many vectors there are nor the order they come out in --- while the branch
\& bound inside each probe grows by a factor of two to six.  A zero objective
gives it nothing to prune with: every node's relaxation is worth $0$, so no
bound can discard a subtree until an integer point has been stumbled upon.

So the direction is load-bearing rather than decorative, the criterion-space
method uses the same one, and the margins stand.  This is reported at length
because it is a plausible-looking optimisation that makes things much worse,
and because it was checked while \emph{suspecting the comparison was unfair} ---
which is the only time such a check is worth anything.

\section{Four improvements to the exact search}

Each of the four removes work without changing the answer.  They are listed
with the unit each was measured in, because no single unit sees all four.

\begin{table}[h]
\centering
\caption{What each change removes, and what it leaves alone.}
\begin{tabular}{lllr}
\toprule
change & unit & before $\to$ after & \\
\midrule
early exit on a hopeless bound & sub-problems & $733 \to 549$ & $-25\%$ \\
not re-deriving Ecker \& Kouada & b\&b nodes & $8105 \to 7618$ & $-6.0\%$ \\
warm-starting a box from its parent & simplex pivots & $18193 \to 16169$ & $-11.1\%$ \\
cutting on a remembered efficient point & simplex pivots & $16169 \to 14602$ & $-9.7\%$ \\
\bottomrule
\end{tabular}
\end{table}

\begin{center}
\fig{r-improve}\\[2pt]
\fig{r-improve-legend}
\end{center}

\subsection{Why four different units}

Because each change removes a different \emph{kind} of work, and a unit that
sees one is blind to the others.

\begin{itemize}
\item The early exit removes whole sub-problems, so it shows in a box count ---
  but those sub-problems are the \emph{cheapest} ones, each killed by the root
  relaxation the moment it sees the cutoff, so it barely moves nodes or pivots.
\item Ecker \& Kouada removes an integer program per inefficient box, so it
  shows in programs and nodes --- but not in the box count, which is identical.
\item The warm start removes a simplex phase~I per box, so it shows only in
  pivots: programs, nodes and boxes are all identical.
\item The remembered-point shortcut removes an efficiency test \emph{and}
  changes which boxes are produced, so it is the only one of the four that is
  not answer-identical by construction and had to be checked against the
  answer instead.
\end{itemize}

\subsection{The measurement lesson}

The wall-clock A/B for the Ecker \& Kouada change reads $0.96\times$ ---
\emph{slower} --- and is faster on 9 of 18 instances.  That cannot be causal:
removing 125 integer programs cannot cost time.

The conclusion is about the instrument.  This machine's noise floor is roughly
$\pm 30\%$ per instance \emph{even at per-instance minima over five repeats},
so it cannot resolve a $6\%$ effect, and no number of repeats fixes that.
Counting integer programs, nodes and pivots gives the same figure on every run
and every machine.

Two habits follow, both adopted only after failing at them first:

\begin{enumerate}
\item A measurement taken while anything else was running on the machine is
  \textbf{discarded, not quoted}.  The first A/B for that change was
  contaminated exactly so and was re-run.
\item A gain predicted from the \textbf{mean} cost of a call is not a
  prediction.  The profile showed the efficiency test at $18.4$\,ms per call,
  the dearest operation in the search --- but what the shortcut removes is the
  repair's test at $7.68$\,ms, the cheap kind.  Reasoning from the wrong row of
  the right profile overstated the gain by half.
\end{enumerate}

\section{Hybridisation: what paid and what did not}

\subsection{Seeding the exact search}

\begin{table}[h]
\centering
\caption{Seeding the box search, 18 instances.  ``Exactly optimal'' counts the
instances where the seed already attains $\Phi^{*}$.}
\begin{tabular}{lrrr}
\toprule
seed & total & seed cost & exactly optimal \\
\midrule
none                   & $9.64$ s  & ---      & --- \\
Pareto local search    & $\mathbf{7.97}$ s $(1.21\times)$ & $0.57$ s & 12 of 18 \\
augmented Tchebychev   & $13.50$ s $(0.71\times)$ & $4.25$ s & 5 of 18 \\
\bottomrule
\end{tabular}
\end{table}

\begin{center}\fig{r-seed}\end{center}

\noindent
The Tchebychev seed is \emph{slower than not seeding at all}, and both halves
of that are structural.  Its cost is $p+1$ integer programs on an enlarged
model --- not the reference point's fault, since computing $z^{*}$ took
$0.01$--$0.13$\,s of the seed against $0.13$--$0.85$\,s for the scalarisations.
Its quality is the real finding: the program steers by criterion-space geometry
relative to $z^{*}$ and \textbf{never looks at $\Phi$}, while the local search
is guided by $\Phi$ at every move.

\paragraph{Thirty paired instances, and one verdict softens.}  Eighteen
instances quoted as a total has the same flaw as the margin table above, so the
campaign re-ran it paired: each of thirty instances at $n=6$, $p=4$ is solved
three times, and the statistic is the median of the per-instance speed-ups.

\begin{table}[h]
\centering
\caption{What a seed is worth, thirty paired instances.  The gap is
$(\Phi^{*}-\Phi(\text{seed}))/|\Phi^{*}|$, zero exactly when the seed is
optimal.}
\begin{tabular}{lrrrr}
\toprule
seed & speed-up (median) & IQR & seed gap & exactly optimal \\
\midrule
Pareto archive & $1.10\times$ & $[1.02,\ 1.27]$ & $0.000$ & $24/30$ \\
Tchebychev     & $1.02\times$ & $[1.00,\ 1.09]$ & $0.170$ & $8/30$ \\
\bottomrule
\end{tabular}
\end{table}

\noindent
The \emph{ordering} is what every sample has said and is the part that holds:
the Pareto seed is worth something, the Tchebychev seed is not.  Two things
change, and both are corrections.

The Pareto seed is worth \textbf{less} than the $1.21\times$ above: a median of
$1.10\times$, with a first quartile of $1.02\times$, so half the instances gain
three percent or less.  The win is real --- both quartiles sit above $1$ and
the seed is exactly optimal on $24$ of $30$ --- but it is small, and the worst
instance runs at $0.618\times$, a genuine loss, because the verification is
paid whether or not the bound can use it.

And the Tchebychev seed \textbf{does not lose here; it does nothing}: not
$0.71\times$ but $1.02\times$, IQR $[1.00,\ 1.09]$ --- a wash.  The verdict
that survives both samples is the weaker one, \emph{it never pays}.  What is
stable is the quality figure that explains it: a median gap of $0.170$ against
$0.000$, and $8$ of $30$ exactly optimal against $24$ of $30$.

\subsection{But as a generator it is the stronger of the two}

A weighted sum maximises only at \emph{supported} efficient points.  Varying
the Tchebychev weights reaches unsupported ones too, which is what the
literature uses it for.

\begin{center}
\fig{r-reach}\\[2pt]
\fig{r-reach-legend}
\end{center}

\begin{table}[h]
\centering
\caption{Reach from the same grid of 27 weight vectors, 20 random instances.}
\begin{tabular}{lrr}
\toprule
 & efficient points reached & of which unsupported \\
\midrule
weighted sum         & $66 / 175$  & $\mathbf{0 / 54}$ \\
augmented Tchebychev & $112 / 175$ & $29 / 54$ \\
\bottomrule
\end{tabular}
\end{table}

\noindent
All 540 optima returned were efficient.  A wider sweep over 40 instances
agrees: $216/369$ against $135/369$, 56 unsupported against 0, and 1080 optima
with none inefficient.

\section{The asymmetry, which is the most transferable result}

\begin{table}[h]
\centering
\caption{The \emph{same} free optimum, handed to each exact method at
iteration 1.}
\begin{tabular}{lrl}
\toprule
exact method & effect & mechanism \\
\midrule
criterion-space box search & $1.42\times$, up to 52\% fewer boxes
  & boxes discarded whole \\
decision-space cut method  & \textbf{zero} iterations saved
  & none \\
\bottomrule
\end{tabular}
\end{table}

\noindent
The heuristic is identical in both rows.  What differs is whether the exact
method has a mechanism through which a lower bound can remove work: a box is
compared against the incumbent \emph{before} it is split, and failing that
comparison deletes it with every descendant it would have had; what the
decision-space method cuts is decided by the efficiency test on its own
maximiser, and $\Phi_{\mathrm{opt}}$ enters only as a cutoff inside a branch
and bound that would have found the same point anyway.

So ``hybridise an exact method with a metaheuristic'' is not one question with
one answer, and the mechanism is cheaper to check than the hybrid is to build
and measure.

\subsection{Two ideas that feed the same mechanism do not add}

\begin{table}[h]
\centering
\caption{Sub-problems solved over the 18 instances.}
\begin{tabular}{lrr}
\toprule
 & boxes solved & reduction \\
\midrule
neither                       & 733 & --- \\
early exit only               & 549 & 25\% \\
metaheuristic seed only       & 562 & 23\% \\
both                          & 539 & 26\% \\
\bottomrule
\end{tabular}
\end{table}

\noindent
Either alone removes about a quarter; together they remove about a quarter.  On
top of the seed the exit is worth $4\%$ rather than $25\%$, because a good
incumbent arrives early enough that the hopeless tail barely forms.  Two ideas
add only when they feed different mechanisms; two that feed the same one are
one idea, measured twice.

\section{Enumerating the whole front, and what the margin tracks}

The same substitution works on complete enumeration of the non-dominated set,
and there the mechanism is visible directly: \textbf{every vector is a cut}, so
in decision space a larger front means more binaries and more big-$M$ rows on
every sub-problem that follows, while the box list never grows.  The campaign
swept the constraint tightness on purpose to move $|F|$ --- fifteen instances
at each of three tightnesses, $n=5$, $p=3$, nothing censored, the same vectors
returned on both sides every time.

\begin{table}[h]
\centering
\caption{The front-enumeration margin against the size of the front, forty-five
instances.}
\begin{tabular}{lrrrr}
\toprule
$|F|$ & instances & ratio (median) & IQR & max \\
\midrule
$1$--$4$  & $26$ & $2.7\times$  & $[1.7,\ 3.9]$   & $9.8\times$ \\
$5$--$12$ & $17$ & $7.5\times$  & $[5.2,\ 16.0]$  & $23.8\times$ \\
$13+$     & $2$  & $31.4\times$ & $[24.9,\ 38.0]$ & $38.0\times$ \\
\midrule
all       & $45$ & $4.1\times$  & $[2.3,\ 7.6]$   & $38.0\times$ \\
\bottomrule
\end{tabular}
\end{table}

\begin{center}
\fig{r-campaign-front}
\end{center}

\noindent
The margin roughly triples and then quadruples as the front grows: $|F|$
\emph{is} the number of cuts.  The $13+$ row rests on two instances and is
quoted as such --- this family at $n=5$ tops out near $|F|=13$.  A separately
measured instance with $|F| = 30$ reached $204\times$; that is an observed
extreme, and the typical case on this family is the $4.1\times$ of the last
row.

\paragraph{One ratio goes the other way.}  Counting branch \& bound nodes,
criterion space explores about three and a half times \emph{more} of them and
still finishes four times sooner --- its nodes carry no binary and no big-$M$
row, so it wins by solving many easy programs instead of few hard ones.  The
same count in the margin table goes the other way at every $p$, because there
the two methods solve nearly the same number of sub-programs and the
decision-space ones are the harder.  Both readings say the same thing: the cost
lives in the size of a sub-problem.

\section{The complete efficient set, for under one doubling}

The front enumeration proves its \emph{vector} list complete and says, in the
same breath, that the \emph{point} list is not: one point is kept per slice.
Closing that gap turns out to be nearly free, for one reason:

\begin{quote}
If $v$ is non-dominated, \textbf{every} $x \in \Dset$ with $\Z(x) = v$ is
efficient --- a $y$ dominating such an $x$ would give $\Z(y) \geqslant v$ with
a strict inequality, making $v$ dominated.
\end{quote}

\noindent
So $\E$ is exactly the union of the front's slices, and \textbf{no efficiency
test is paid on any slice point}: the dearest operation in the package is not
invoked at all.  A slice is $\Dset$ plus $p$ linear equalities --- linear even
for fractional criteria, since $d^{\top}x + b > 0$ clears the ratio --- so
there is no binary, no big-$M$, and no growth from one slice to the next.

Below, five variables enter $\Z$ and $f$ further ones enter only their own
bound $0 \leqslant x_j \leqslant 4$: invisible to $\Z$, so each slice is a grid
of $5^{f}$ points that the front collapses to one.

\begin{table}[h]
\centering
\caption{What the slices recover, and what they cost.}
\begin{tabular}{rrrrrr}
\toprule
$f$ & $|$front$|$ & $|\E|$ & ratio & slices & extra time \\
\midrule
0 &  7 &     7 & $1\times$   & $0.016$\,s & $24\%$ \\
0 & 13 &    13 & $1\times$   & $0.018$\,s & $10\%$ \\
1 &  7 &    35 & $5\times$   & $0.019$\,s & $27\%$ \\
1 & 13 &    65 & $5\times$   & $0.024$\,s & $13\%$ \\
2 &  7 &   175 & $25\times$  & $0.036$\,s & $42\%$ \\
2 & 13 &   325 & $25\times$  & $0.050$\,s & $24\%$ \\
3 &  7 &   875 & $125\times$ & $0.103$\,s & $90\%$ \\
3 & 13 &  1625 & $125\times$ & $0.167$\,s & $68\%$ \\
\bottomrule
\end{tabular}
\end{table}

\begin{center}
\fig{r-complete-gain}\\[6pt]
\fig{r-complete-cost}\\[2pt]
\fig{r-complete-legend}
\end{center}

\noindent
The two charts are the finding: what is recovered multiplies by five at each
step, and the extra time does not.  At $f = 3$ the front is \emph{complete and
missing $99.2\%$ of the points}, and recovering all $1625$ of them never
doubles the cost.  The scan is output-sensitive --- about $0.1$\,ms per point,
against $7.68$\,ms for a single efficiency test, of which none are paid.

\paragraph{The honest reading, which is conditional.}  At $f = 0$ the front
\emph{already was} the efficient set, and the $10$--$24\%$ buys only the proof
of that --- the usual case on generic instances, where $476$ efficient points
sat on $476$ distinct vectors over $32$ of them.  This is not a speed-up and it
is not uniform.  What it buys is that the method now \emph{knows} which case it
is in, rather than returning a point list whose completeness rested on an
assumption about the coefficients.

\section{The hybridisation that pays, and the three that did not}

Nothing above was both \emph{hybrid} and \emph{complete}: $\E$ was reached by a
purely exact method, and every hybridisation had been hung on the optimisation
search, where a seed feeds the \emph{bound}.  The front enumeration has no
bound and no incumbent --- it must exhaust every box --- so a seed of that kind
is worth nothing there.

It has a different mechanism.  Every box costs a probe, an efficiency test and
sometimes a repair, all spent to find \emph{a vector}; a vector already known
needs none of them.  Seeding here is not ``start from a better value'' but
``do not pay to discover what you already know''.

\subsection{First attempt: a generator.  It loses.}

The augmented Tchebychev program is the natural candidate --- its optima are
efficient by construction, so they need no verification.  Over 120 instances it
is \emph{worse the more of the front it seeds}: with a wide weight grid it
covers $100\%$ of the front and runs at $2.61\times$ the unseeded time.

The reason is one measurement: a scalarisation is not cheaper than the probe it
replaces.

\begin{table}[h]
\centering
\caption{One probe against one scalarisation, identical instances.}
\begin{tabular}{lrr}
\toprule
 & time & b\&b nodes \\
\midrule
one probe         & $9.48$\,ms  & $0$ \\
one scalarisation & $20.90$\,ms & $21$ \\
\bottomrule
\end{tabular}
\end{table}

\noindent
Cheaper on only 7 of 40.  Seeding $k$ vectors saves at most $k$ probes and
costs at least $k$ scalarisations plus $p$ for the ideal point --- a losing
trade before the duplicate weights are counted.

\subsection{Second attempt: a walk.  It pays --- once it is sorted.}

A Pareto local search is different in kind: one walk of about $0.011$\,s yields
a whole archive covering $95\%$ of the front.  Its points need verification,
and that verification is \textbf{not} a cost here --- the enumeration was going
to pay the same efficiency test on every probe the seed removes.  It cancels,
and the probe is the gain.

But handed over in archive order the hybrid \emph{lost}, at $0.81\times$, with
the probe count \emph{rising} on 4 of 12 instances: the pre-split was building
a box structure the search would not have chosen.  Sorting the seeds by the
direction the probe itself maximises reproduces that structure instead.

\begin{center}
\fig{r-order}
\end{center}

\begin{table}[h]
\centering
\caption{Probes against the unseeded search, by the order the seeds arrive.}
\begin{tabular}{lrr}
\toprule
order & probes & raised the count on \\
\midrule
archive order        & $0.89\times$ & 4 of 12 \\
\textbf{probe order} & $\mathbf{0.70\times}$ & \textbf{none} \\
reversed             & $0.94\times$ & 5 of 12 \\
\bottomrule
\end{tabular}
\end{table}

\noindent
The reversed arm makes the direction causal rather than a coincidence.
\textbf{A $1.33\times$ win was sitting behind a sort.}

\begin{table}[h]
\centering
\caption{The hybrid against the pure exact methods.  Thirty instances at each
of five configurations, $n = 6 \dots 9$; medians of per-instance ratios.}
\begin{tabular}{lrrr}
\toprule
 & median & IQR & faster on \\
\midrule
front enumeration & $1.33\times$ & $[1.20,\ 1.51]$ & $139/150$ \\
$\E$ complete     & $\mathbf{1.31\times}$ & $[1.17,\ 1.47]$ & $139/150$ \\
probes            & $0.72\times$ & $[0.67,\ 0.75]$ & \emph{never worse} \\
\bottomrule
\end{tabular}
\end{table}

\begin{center}
\fig{r-hybrid}
\end{center}

\noindent
Identical answers on all 150 --- same front, same $\E$, same optimum --- which
the campaign asserts rather than assumes.  The heuristic decides only how much
of the front has to be discovered.

\section{Solving the slice equations: the largest gain here}

Completing the front to $\E$ was dominated by the slice scan: at $n \geqslant 9$
it took $25$--$30$\,s against under a second for the front, and at $n = 12$ a
slice could not finish inside two minutes.

The $p$ slice equations are \emph{linear}, so solve them.  Gaussian elimination
in exact rationals expresses $r$ variables as functions of the rest;
substituting them into every model row leaves an equivalent system in $n - r$
unknowns.  The search space $(ub+1)^{n}$ becomes $(ub+1)^{n-r}$.

\begin{table}[h]
\centering
\caption{Pruning on the slice equations against solving them.  Fourteen
instances, $n = 7 \dots 12$, the same points every time.}
\begin{tabular}{lrrr}
\toprule
instance & prune & eliminate & ratio \\
\midrule
$n=9$, $p=3$  & $29.71$\,s  & $0.31$\,s & $96.7\times$ \\
$n=10$, $p=5$ & $119.47$\,s & $0.92$\,s & $129.9\times$ \\
$n=12$, $p=3$ & $796.84$\,s & $8.34$\,s & $95.5\times$ \\
\midrule
median of all fourteen & & & $\mathbf{49.2\times}$ \\
\bottomrule
\end{tabular}
\end{table}

\paragraph{A sign error that ran, and returned nothing.}  The first draft put a
plus where the eliminated variables' contribution must be subtracted.  It ran
without error, looked like a working optimisation, and returned the
\emph{empty} set on 179 of 185 slices --- silently dropping efficient points.
Only the comparison against the un-eliminated scan caught it.

\section{A wall met four times}

Four separate changes remove between a fifth and a quarter of the sub-problems
and buy nothing.

\begin{table}[h]
\centering
\caption{Four attempts on the sub-problem count.}
\begin{tabular}{lrr}
\toprule
change & sub-problems removed & cost \\
\midrule
early exit on a hopeless bound    & $25\%$       & no measurable time \\
free contradiction filter         & $24.8\%$     & $0.98\times$ pivots \\
range filter                      & $+11.3\%$    & $1.03\times$ pivots \\
dominated-whole filter            & $8$--$23\%$  & $1.077\times$ pivots \\
\bottomrule
\end{tabular}
\end{table}

\begin{center}
\fig{r-wall}\\[2pt]
\fig{r-wall-legend}
\end{center}
\noindent
\emph{Three of the four cost time rather than saving it; a bar cannot be drawn
negative, so they sit at zero.  The gap between the two series is the finding.}

\begin{quote}
\textbf{In this search, removing sub-problems does not remove time, because the
ones that can be removed cheaply are the ones that were already cheap.}
\end{quote}

\noindent
An empty box --- and $73\%$ of them are --- is refuted by the simplex almost at
once: every probe in these families costs \emph{zero} branch \& bound nodes,
because the relaxation is integral or infeasible at the root.  A filter that
catches empty boxes catches the cheapest work in the search.  The only change
that went round this wall attacked an \emph{exponent} rather than a count, and
is the $49\times$ above.

\paragraph{The same wall on the proof.}  After the walk seeds $95\%$ of the
front, $93\%$ of the remaining boxes are empty and the enumeration is almost
purely a \emph{proof of emptiness}.  ``Is there a feasible point not dominated
by any of the $k$ known ones?'' is \textbf{one} question, and Sylva--Crema
answers it in one integer program built once, with no iteration and no growth.
It is sound --- 13 of 30 instances proved complete by a single program, no
wrong verdict --- and $4.5\times$ slower in the median, $100\times$ at worst
($239$\,s against $1.95$\,s), because $k$ cuts mean $kp$ binaries in one model.

That answers a question wider than the optimisation.  Criterion space was
claimed to be the better \emph{search} because the model never grows; this
tests whether it is also the better \emph{proof}, where nothing grows at all.
\textbf{It is, by $4.5\times$.}

\section{Fractional criteria, where the problem lives}

Every figure so far uses \emph{linear} criteria --- a gap worth naming in a
package about ratios.  It has a reason: the Tchebychev comparison needs them,
since $z^{*}_k - \Z_k(x)$ is not linear when $\Z_k$ is a ratio.  The Pareto
walk has no such requirement.

\begin{table}[h]
\centering
\caption{Genuinely fractional criteria, ten instances per row.}
\begin{tabular}{lrrr}
\toprule
configuration & $|F|$ & cover & ratio \\
\midrule
$n=5$, $p=2$ & $7$  & $100\%$ & $1.54\times$ \\
$n=6$, $p=3$ & $22$ & $100\%$ & $1.35\times$ \\
$n=7$, $p=3$ & $36$ & $100\%$ & $1.25\times$ \\
\midrule
all $40$ & & & $\mathbf{1.39\times}$, faster on $34$ \\
\bottomrule
\end{tabular}
\end{table}

\noindent
At least as good as the linear case, with the front covered completely.
Nothing in the construction distinguishes the two.

\section{Five families, and a prediction that was wrong}

Everything else here comes from one generator, so the magnitudes could belong
to it.  Four more families move what the mechanism depends on: \texttt{loose}
raises $|\Dset|$ an order of magnitude while leaving $|F|$ alone,
\texttt{knapsack} is two capacity rows with all coefficients positive,
\texttt{conflicting} gives the criteria opposing signs (large $|F|$), and
\texttt{aligned} correlates them (small $|F|$).

\begin{table}[h]
\centering
\caption{Fifteen instances per cell, 150 in all; two sizes per family.
Identical front and identical optimum on every one.}
\begin{tabular}{lrrrr}
\toprule
family & $|F|$ & cover & ratio & faster on \\
\midrule
\texttt{mixed}       & $8$ / $30$    & $100$ / $97\%$ & $\mathbf{1.37}$ / $1.30\times$ & $28/30$ \\
\texttt{loose}       & $9$ / $44$    & $83$ / $86\%$  & $1.29$ / $1.23\times$ & $28/30$ \\
\texttt{knapsack}    & $12$ / $41$   & $100$ / $90\%$ & $1.33$ / $1.27\times$ & $29/30$ \\
\texttt{conflicting} & $107$ / $105$ & $100$ / $93\%$ & $1.14$ / $1.14\times$ & $26/30$ \\
\texttt{aligned}     & $6$ / $11$    & $100$ / $84\%$ & $1.31$ / $1.14\times$ & $25/30$ \\
\midrule
\textbf{all 150} & & & $\mathbf{1.23\times}$ & $\mathbf{136}$ \\
\bottomrule
\end{tabular}
\end{table}

\begin{center}
\fig{r-families}\\[2pt]
\fig{r-families-legend}
\end{center}

\noindent
Every family pays, $1.14$--$1.37\times$.  The advantage is not an artefact of
the generator it was found with, which is what this was run for.

\paragraph{And the prediction it refutes.}  We expected \texttt{conflicting},
with $|F| \approx 107$, to be the family the hybrid liked \emph{most} --- a
large front is much to discover, and discovery is what a seed removes --- and
\texttt{aligned} with $|F| = 6$ to be the least.  \textbf{The opposite holds},
and in the chart the orange bar is below the blue one in four families and
level with it in the fifth --- never above: the margin narrows as $|F|$ grows,
within families as well as across them.

\paragraph{The ceiling that explains it.}  A seed removes \emph{one} probe, and
the enumeration spends $2.6$ to $4.5$ probes per vector --- the rest go to
boxes that turn out empty, which no seed can touch.  So even at full coverage
the probe saving is capped near a quarter to a third, exactly the $24$--$32\%$
the campaign measures.

\begin{quote}
\textbf{That is a hard ceiling of roughly $1.4\times$ on the whole idea, and it
does not depend on $|F|$.}  What does depend on $|F|$ is the walk's cost, one
efficiency test per archive member --- so a large front buys no more and costs
more.
\end{quote}

\section{Where the hybrid fails, and the line that fixes it}

At $n = 12$ the margin falls to $1.00\times$, and at $n=12$, $p=5$ to
$0.97\times$.  Our first explanation was that the probe count outgrows $|F|$.
\textbf{The measurement refutes it}: probes per vector are flat across the
range --- $2.4$, $2.7$, $2.7$, $2.8$, $2.8$, $3.0$, $3.9$, $5.2$.  What
collapses is the share of the front the walk finds, and the speed-up tracks it.

\begin{center}
\fig{r-cover}\\[2pt]
\fig{r-cover-legend}
\end{center}

\begin{table}[h]
\centering
\caption{Widening the walk alone, on the two instances that failed.}
\begin{tabular}{lrrrr}
\toprule
 & \multicolumn{2}{c}{$n=12$, $p=3$} & \multicolumn{2}{c}{$n=12$, $p=5$} \\
restarts / budget & cover & ratio & cover & ratio \\
\midrule
$8$ / $4000$    & $35\%$ & $1.03\times$ & $12\%$ & $0.98\times$ \\
$16$ / $20000$  & $95\%$ & $\mathbf{1.61\times}$ & $46\%$ & $\mathbf{1.20\times}$ \\
$32$ / $60000$  & $95\%$ & $1.53\times$ & $93\%$ & $1.12\times$ \\
$64$ / $150000$ & $95\%$ & $1.44\times$ & $89\%$ & $1.17\times$ \\
\bottomrule
\end{tabular}
\end{table}

\noindent
Restarts are the knob, not neighbours, and the fix is one line: the first round
is sized to the instance.  \emph{Not} a second constant --- $4000$ was chosen
at $n=7,8$ and generalised, which is what produced the collapse.

\paragraph{But coverage is not the objective.}  The $p=5$ column says so
directly: \textbf{$46\%$ coverage gives $1.20\times$ and $93\%$ gives
$1.12\times$}.  A walk that restarts until it finds nothing new takes coverage
there from $45\%$ to $97\%$ and is \emph{slower} --- $1.02\times$ against
$1.14\times$, and a median $1.32\times$ against $1.36\times$ over 120
instances.  It is in the code and switched off.

\paragraph{A correction to our own reading.}  The collapse was first reported
as ``$35\%$ coverage, $1.00\times$'' from \emph{one} instance; over six seeds
of that configuration the unmodified walk averages $78\%$ and $1.35\times$.
Real and severe where $|F|$ is large, but not as general as one instance made
it look --- the fifth time here that a single instance misled.

\section{How the campaign was measured}

Most figures here were first taken on four to eight instances.  The
\emph{direction} of each was argued from a mechanism, which is what made them
worth reporting --- but a total over four runs is not an estimate.
\texttt{studies/campaign.py} re-measures four of the claims over thirty
instances per configuration, in $751$ seconds, and writes every per-instance
record to a JSON file so the tables can be rebuilt and re-checked without
paying for the run again.

\begin{itemize}
\item \textbf{The statistic is the median of the per-instance ratios}, with the
  interquartile range, never a ratio of totals.  A ratio of totals is decided
  by the single heaviest instance --- which is exactly how the $103\times$ at
  $p=5$ came about.
\item \textbf{Deterministic work is counted next to the seconds}: sub-programs
  solved and branch \& bound nodes explored, identical on every run and every
  machine.  This machine's noise floor is roughly $\pm 30\%$ per instance, so
  seconds alone cannot resolve a small effect and no number of repeats fixes
  that.
\item \textbf{Censoring is reported, never hidden.}  A run that spends its time
  budget is excluded from the medians and counted in the open, and a row with
  more than half its runs censored says so instead of quoting a number.  In the
  margin experiment censoring is one-sided --- it can only flatter the
  decision-space method --- so those margins are conservative.
\item \textbf{Agreement is asserted, not assumed.}  Every experiment checks
  that the two methods return the same answer, not merely that one is faster.
\end{itemize}

\noindent
Nothing else ran on the machine during the campaign; a contaminated measurement
is discarded rather than quoted, which is a rule adopted here only after
breaking it once.

\section{Every verdict, including the losses}

\begin{table}[h]
\centering
\caption{The twenty-two things tried, with what each was worth.}
\begin{tabular}{lll}
\toprule
attempt & verdict & effect \\
\midrule
decision space $\to$ criterion space & \textbf{win}  & $1.5\times$ at $p=2$ to $14.6\times$ at $p=6$ \\
exact--exact handover ($s=1$)        & tie           & $8.47$ s vs $8.11$ s \\
exact--exact handover ($s=3$)        & \textbf{loss} & $13.14$ s vs $8.47$ s \\
exact--metaheuristic (Pareto archive)& \textbf{win}  & $1.10\times$ median, $24/30$ seeds exact \\
exact--Tchebychev seed               & \textbf{loss} & $1.02\times$: a wash, never pays \\
Tchebychev as a generator            & \textbf{win}  & reaches unsupported points \\
incumbent seeding in decision space  & \textbf{loss} & zero iterations saved \\
LP pre-filter on each box            & \textbf{loss} & $7.95 \to 8.36$ s, box counts identical \\
two-at-a-time split                  & \textbf{loss} & loses across $p=3..7$, worse as $p$ grows \\
early exit on a hopeless bound       & partial       & 25\% fewer boxes, no measurable time \\
early exit \emph{on top of} the seed & \textbf{loss} & 4\% rather than 25\%: a substitute \\
warm start + remembered point        & \textbf{win}  & 20\% fewer pivots together \\
zero-objective probe                 & \textbf{loss} & $4$--$13\times$ slower, cuts identical \\
slices instead of an efficiency test  & \textbf{win}  & all of $\E$, under $2\times$ \\
eliminating the slice variables       & \textbf{win}  & $49\times$ median, up to $130\times$ \\
Tchebychev seeding the enumeration    & \textbf{loss} & worse the more it seeds \\
\textbf{Pareto walk seeding it}              & \textbf{win}  & $\mathbf{1.31\times}$ to all of $\E$ \\
free box filters (three of them)      & partial       & $2\%$, then two losses \\
one-program completeness certificate  & \textbf{loss} & $4.5\times$ slower, sound \\
sizing the walk to the instance       & \textbf{win}  & $1.00 \to 1.53\times$ at $n=12$ \\
a walk that restarts until it stops   & \textbf{loss} & more coverage, less speed \\
the hybrid on fractional criteria     & \textbf{win}  & $1.39\times$, better than linear \\
\bottomrule
\end{tabular}
\end{table}

\noindent
Nine are losses and one is a tie.  They are reported at the same length as the
wins because they carry the same information: a change pays only where the
method has a mechanism it can feed, that mechanism is cheaper to identify than
the change is to build and measure, and two ideas feeding the same mechanism
are one idea counted twice.

\section{How all of this is checked}

Exact rational arithmetic throughout, no floating point anywhere, no
third-party dependency, and 93 tests run on every commit.  Several check
\emph{properties} rather than answers, since a method that silently lost
efficient points would still agree with the optimum on favourable instances:
the split is a partition; the metaheuristic never hands over an inefficient
point; every remembered point is efficient; a warm root reaches the same value
as the cold one it replaces; the early exit leaves boxes open and still proves
optimality.

Answers are cross-checked three ways: against the decision-space method,
against exhaustive enumeration on small instances, and against a
$\Phi$-ordered scan on instances up to $34635$ feasible points.

\end{document}
